\section{Online Learning：可能的范式级路线}

前面几章说明 Agent 为什么需要环境和反馈，本章进入 Online Learning。广密认为，如果未来还有范式级新路线，Online Learning 是候选之一。它的核心不是训练完模型再部署，而是让模型在环境中自主探索、获得反馈、更新经验，甚至进一步更新权重或长期记忆。

\begin{figure}[H]
\centering
\includegraphics[width=0.96\textwidth]{online-learning-loop.png}
\caption{Online Learning 闭环：模型在环境里自主探索、获得反馈、更新能力。自制概念图，依据 00:55:49--01:02:45 对谈内容整理。}
\end{figure}

\begin{knowledgebox}{读图：Online Learning 把数据生产放进使用过程}
从 Environment 到 Action、Reward、Update、Memory/Weights，图里强调的是持续循环。传统训练更像离线学习；Online Learning 则让模型在真实或仿真环境中边做边学。Agent 越多，潜在反馈数据越多。
\end{knowledgebox}

\subsection{为什么 Online Learning 可能是下一范式}

本节说明其吸引力。Pre-training 主要依赖已有语料，Post-training/RL 依赖整理好的任务和奖励；Online Learning 则可能从模型自己的行动中产生新数据。对于 Coding、游戏、科学实验、网页操作等环境，模型可以尝试、失败、修复、总结，再把经验回流。它把数据瓶颈从“找更多旧数据”转向“创造更高价值的新经验”。

\begin{importantbox}{核心转变：从数据集到环境}
如果 Online Learning 成立，训练中心会从静态数据集转向可交互环境。谁拥有环境、任务、奖励、评测和真实用户反馈，谁就可能拥有新的数据飞轮。
\end{importantbox}

\subsection{对 GPU 和 NVIDIA 叙事的影响}

本节回应访谈中的产业追问。Online Learning 不一定削弱 GPU 需求，反而可能改变 GPU 需求结构：不仅要训练，还要大规模采样、执行、多轮推理、评测和回放。推理和训练可能更紧密交织，infra 需要同时支持在线 rollout、缓存、环境模拟和安全审计。NVIDIA 的叙事也不只是卖训练卡，而是围绕推理、仿真、机器人、数据中心和软件栈扩展。

\begin{knowledgebox}{术语消化：Online Learning 相关基础设施}
\noindent\begin{tabularx}{\textwidth}{p{0.20\textwidth}p{0.34\textwidth}X}
\toprule
术语 & 含义 & 为什么重要 \\
\midrule
Rollout & Agent 在环境中执行一段轨迹 & 产生可评价的新经验。 \\
Reward & 评价任务成功与质量的信号 & 决定模型会学到什么。 \\
Replay & 回放成功/失败轨迹用于训练或分析 & 帮助从错误中学习。 \\
Simulation & 可控环境或仿真世界 & 降低真实试错成本。 \\
Memory Update & 把经验写入长期记忆或外部状态 & 让模型越用越懂。 \\
\bottomrule
\end{tabularx}
\end{knowledgebox}

\subsection{风险：在线学习如何不被 reward hack}

本节补充约束。Online Learning 的风险在于，模型可能学会钻奖励、污染记忆、放大偏见、绕过安全边界，或者从错误反馈中学到错误策略。因此真正可用的在线学习系统必须有分层权限、沙盒环境、可回滚状态、离线评测、人工审计和安全过滤。

\begin{warningbox}{Online Learning 不能等同于“让模型自己随便学”}
持续学习如果没有边界，会带来记忆污染、目标漂移和安全问题。高质量 Online Learning 更像受控实验系统：任务清楚、奖励可审、失败可追踪、更新可回滚。
\end{warningbox}

\subsection{本章小结}

Online Learning 的想象力在于，它把 Agent 的使用过程变成数据生产过程。它也带来新的基础设施和安全问题。下一阶段，谁能稳定运行环境、奖励和反馈闭环，谁就可能获得新的模型改进路径。

